\chapter{Podsumowanie i~wnioski końcowe}
\label{c:podsumowanie}

Po zakończeniu badań i~eksperymentów przeprowadzonych w~ramach niniejszej pracy można sformułować kilka istotnych wniosków i~refleksji dotyczących skuteczności zastosowania metod detekcji opartych na uczeniu głębokim w~kontekście rozpoznawania obiektów trzymanych przez osoby niewidome i~słabowidzące. Praca objęła część projektowo-implementacyjną oraz badawczą. Ten rozdział stanowi podsumowanie wykonanych prac, ich ocenę oraz wskazanie potencjalnych kierunków dalszego rozwoju systemu.

\section{Podsumowanie wykonanych prac}

Pierwszym etapem pracy było poszerzone studium literatury, które uwidoczniło kierunki rozwoju technologii wspomagających osoby niewidome i~słabowidzące. Pośród nich znajdujemy wiele rozwiązań opartych na przetwarzaniu obrazu, które umożliwiają wsparcie w~codziennych czynnościach. Nie zidentyfikowano w~nim jednak żadnej pracy, która opisywałaby wsparcie osób niewidomych w~nauce geometrii przestrzennej poprzez interakcję z~fizycznymi modelami trójwymiarowymi. W~związku z~tym, należało opracować własny system, który byłby odpowiedzią na zidentyfikowaną lukę w~literaturze.

Rozwiązanie oparto na architekturze klient-serwer i~metodach detekcji obiektów. Koncepcyjnie rozdzielono odpowiedzialności pomiędzy urządzeniem mobilnym, pełniącym rolę klienta, a~serwerem, który przetwarzał dostarczone informacje i~zwracał wyniki detekcji. Wyznaczono również wymagania funkcjonalne i~niefunkcjonalne systemu, które stanowiły podstawę do dalszych prac.

Założenia, które przyjęto dla systemu, zostały następnie zrealizowane poprzez stworzenie funkcjonalnego prototypu. W~jego ramach powstały aplikacja kliencka przeznaczona na urządzenia mobilne z~systemem Android, serwer przetwarzający dane oraz zestaw modeli detekcji obiektów. Zadana architektura systemu umożliwiła połączenie różnych komponentów w~spójną całość, co sprawiło, że prototyp może zostać uznany za funkcjonalny i~gotowy do dalszego rozwoju. Interakcja z~użytkownikiem końcowym została oparta na komunikacji głosowej, co wykluczyło konieczność korzystania z~ekranów dotykowych i~odrywania rąk od trzymanych obiektów. Stworzono również zestaw narzędzi do przygotowania danych treningowych i~treningu modeli.

Aby możliwe było użycie modeli detekcji, konieczne było przygotowanie odpowiedniego zbioru danych. W~tym celu nagrano i~przetworzono zestaw filmów z~których wyodrębniono obrazy zawierające osoby trzymające w~dłoniach różne bryły geometryczne. Łącznie na zbiór składa się 4666 ręcznie adnotowanych obrazów, które zostały następnie podzielone na zbiory treningowy, walidacyjny i~testowy.

Część badawcza pracy obejmowała eksperymenty z~udziałem siedmiu modeli YOLO26 trenowanych w~różnych konfiguracjach. Celem tej części było zbadanie możliwości ich użycia w~implementowanym systemie. Badano wpływ skali modelu, rozdzielczości wejściowej i~reprezentacji obrazu. Ostatecznie wskazano model, który wykazywał najlepszy kompromis jakości i~kosztu obliczeniowego, charakteryzując się krótkim czasem inferencji przy zachowaniu korzystnej skuteczności detekcji. 

\section{Ocena realizacji celu pracy}

Celem pracy było opracowanie systemu wspomagającego osoby niewidome i~słabowidzące w~nauce geometrii przestrzennej za pomocą fizycznych modeli 3D. Cel ten zrealizowano poprzez stworzenie funkcjonalnego prototypu systemu, który umożliwia rejestrowanie i~przesyłanie obrazu, detekcję brył trzymanych przez użytkownika, analizę widocznej części obiektu oraz przekazywanie informacji w~formie głosowej. Zaimplementowane rozwiązanie realizuje więc podstawowy scenariusz użytkowy przyjęty na etapie projektowania.

Przeprowadzono również eksperymentalną ocenę wytrenowanych modeli detekcji. Uzyskane wyniki potwierdziły, że modele były w~stanie wykrywać wszystkie sześć klas obiektów uwzględnionych w~pracy. Badania umożliwiły wskazanie jednego z~modeli jako najkorzystniejszego kompromisu pomiędzy skutecznością detekcji a~kosztem obliczeniowym i~czasem inferencji. Nie oznacza to jednak, że jest on w~pełni odporny na zmiany środowiska, ponieważ to nie było przedmiotem osobnego eksperymentu i~nie zostało zweryfikowane w~warunkach docelowych.

Zaimplementowany prototyp pozwala na jego działanie w~kontrolowanych warunkach oraz przeprowadzenie dalszych testów z~udziałem osób niewidomych i~słabowidzących. Wdrożenie go w~środowisku docelowym wymaga jednak dalszych badań. W~szczególności należy przeprowadzić ankiety wśród użytkowników końcowych, które pozwolą na ocenę jego użyteczności i~wpływu na proces nauki. Może to wskazać obszary wymagające poprawy i~dalszego rozwoju.

Pełny zakres funkcjonalny aplikacji mobilnej został zrealizowany wyłącznie dla wersji przeznaczonej na system Android. Szkielet implementacyjny dla systemu iOS został przygotowany, jednak nie zawiera ważnych funkcjonalności, co uniemożliwia jego użycie w~przyjętym scenariuszu działania bez dalszych prac. 

Cel pracy został osiągnięty w~przyjętym zakresie. Stworzono działający prototyp systemu i~przygotowano i~oceniono jego kluczowy moduł detekcji obiektów. Jednocześnie zidentyfikowano obszary wymagające dalszych badań i~rozwoju przed jego ewentualnym wdrożeniem w~środowisku docelowym.

\section{Najważniejsze wnioski}

Realizacja pracy pozwoliła na sformułowanie kilku istotnych wniosków, które mogą stanowić podstawę do dalszych badań i~rozwoju systemu.

Wpływ rozdzielczości obrazu wejściowego okazał się jednym z~najistotniejszych czynników wpływających na skuteczność detekcji. Modele trenowane na obrazach o~wyższej rozdzielczości wykazywały ogólną poprawę wyników względem wariantów używających obrazów o~niższej rozdzielczości. Dla modelu YOLO26n poprawie uległy wszystkie badane metryki, natomiast dla modelu YOLO26s poprawa dotyczyła wyłącznie metryk precyzji, $mAP50$ i~$mAP50-95$ i~$F1$-score przy jednoczesnym spadku czułości. Warto jednak zauważyć, że zwiększenie tego parametru wiąże się ściśle ze zwiększeniem wymagań obliczeniowych, prowadząc do wydłużenia czasu inferencji, dlatego jej wybór powinien uwzględnić również kompromis pomiędzy jakością detekcji a~wydajnością systemu. 

Zwiększenie skali modelu nie przyniosło jednoznacznej poprawy skuteczności detekcji. Przy mniejszych badanych obrazach wyższe warianty modeli wykazywały lepsze wyniki, wskazując na korzyści z~zastosowania bardziej złożonych modeli. Dla obrazów o~wyższej rozdzielczości ta zależność nie była już tak wyraźna: występowały wahania poszczególnych metryk, a~w~niektórych przypadkach mniejsze modele osiągały lepsze wyniki. Zmiana skali modelu nie zapewnia więc stałej przewagi, a~zasadność jej stosowania powinna być oceniana w~kontekście innych czynników.

Usunięcie informacji o~kolorze z~obrazów wejściowych poskutkowało znacznym spadkiem skuteczności detekcji badanych modeli. Analiza wyników poszczególnych klas wykazała jednak, że nie było ono jednakowe dla wszystkich klas. Wartości porównywalne z~modelami pracującymi na obrazach kolorowych uzyskano dla klas sześcianu i~prostopadłościanu. Najgorsze wyniki uzyskano dla klasy stożka. Wskazuje to na istotną rolę informacji o~barwie, ale jej znaczenie może być różne w~zależności od klasy obiektu. 

Poszczególne zmiany konfiguracji modeli wpływały w~różnym stopniu na zdolność rozróżniania klas obiektów. Szczególnie często mylone były obiekty o~zbliżonym kształcie, jak sześcian i~prostopadłościan, które w~wielu przypadkach zostawały błędnie przypisane do siebie. 

Uruchomienie detekcji na zbiorach testowych skutkowało osiągnięciem niższych wyników niż w~przypadku zbiorów walidacyjnych. Ukazuje to znaczenie wykonywania testów na danych, które nie uczestniczyły w~procesie treningu. Warto również zauważyć, że dane w~zbiorze testowym były niezależne od zbioru treningowego i~walidacyjnego, przez co stanowiły bardziej niezależny materiał do oceny. Zaobserwowana różnica wskazuje również, że wyniki walidacyjne nie odzwierciedlały w~pełni zdolności modeli do generalizacji.

Podczas wyboru modelu przeznaczonego do zastosowania w~systemie należy zwrócić uwagę nie tylko na jego jakość detekcji, ale również na koszt obliczeniowy i~czas inferencji. Wysoka wartość metryk detekcji nie musi oznaczać, że dany model będzie najlepszym wyborem w~kontekście całego rozwiązania. Model, który został wybrany do implementacji w~systemie, nie osiągał najwyższych wyników, ale wymagał mniejszej liczby zasobów obliczeniowych i~charakteryzował się krótszym czasem inferencji od pozostałych. 

Wnioski, jakie płyną z~przeprowadzonych badań wskazują, że możliwe jest wykorzystanie metod detekcji obiektów w~kontekście tematu pracy. Wskazują również na obszary wymagające dalszych badań i~rozwoju, które mogą przyczynić się do zwiększenia skuteczności, a~w~przyszłości, być może, wdrożenia systemu w~środowisku docelowym.

\section{Możliwe kierunki rozwoju}

W~toku prac zidentyfikowano kilka obszarów, które mogą stanowić kierunki dalszego rozwoju systemu.

Najważniejszym z~nich jest zdecydowanie powiększenie i~zróżnicowanie zbioru danych. Rozmiary obecnego zbioru zostały ograniczone przez czas i~możliwości przeprowadzenia nagrań, a~także konieczność ręcznej adnotacji danych. W~przyszłości warto rozważyć wykorzystanie w~nagraniach większej liczby osób, pomieszczeń i~warunków, co pozwoli na zwiększenie różnorodności materiałów. Warto również rozważyć wprowadzenie dodatkowych klas obiektów, jak np. pozostałych wielościanów foremnych, których własności geometryczne są często omawiane w~ramach nauki geometrii przestrzennej. Zwiększenie liczby klas obiektów może również przyczynić się do zwiększenia użyteczności systemu, poprzez rozszerzenie dostępnych możliwości wsparcia. 

Kolejnym kierunkiem, jaki należy podjąć jest rozszerzenie badań dotyczących wykorzystania modeli pracujących w~odcieniach szarości. Warto przeprowadzić eksperymenty z~dodatkowymi zestawami brył, aby zbadać wpływ ich kolorystyki na skuteczność detekcji. Niewykluczone, że dodanie do zbioru danych brył o~różnych kolorach, mających ten sam kształt, pozwoliłby modelom na lepszą generalizację i~zmniejszenie ich zależności od informacji o~kolorze. 

Następnie, warto wykonać dodatkowe eksperymenty z~wykorzystaniem innych architektur sieci neuronowych, a~także powtórzyć treningi z~wykorzystaniem innych metod augmentacji danych i~wartościami ziarna. Pozwoli to ocenić zmienność wyników i~sprawdzić powtarzalność uzyskanych rezultatów. Warto również uzupełnić zbiór modeli o~wariant YOLO26s pracujący na obrazach w~odcieniach szarości i~rozdzielczości 1024 pikseli. Takie rozszerzenie badań pozwoliłoby dokładniej ocenić wpływ tych czynników na skuteczność detekcji. 

Innym kierunkiem rozwoju jest zmiana sposobu analizy obrazów. Obecna wersja traktuje każdy obraz niezależnie. Zmiana tego podejścia na analizę sekwencji obrazów mogłaby ograniczyć wpływ chwilowego zasłonięcia obiektu, a~także zwiększyć skuteczność detekcji poprzez wykorzystanie informacji zawartej w~kilku kolejnych klatkach. 

Odchodząc od samej detekcji, uwagę zwrócić należy na użytkowników końcowych. Badania przeprowadzone z~osobami niewidomymi i~słabowidzącymi pozwolą na ocenę użyteczności systemu oraz wskażą obszary wymagające poprawy. Zmiany interfejsu użytkownika mogłyby ograniczyć lub wyeliminować konieczność obecności operatora w~trakcie korzystania z~systemu, zwiększając jego samodzielność i~komfort użytkowania. Warto również rozważyć wprowadzenie dodatkowych komend głosowych, które pozwoliłyby na sterowanie całym systemem w~sposób bezdotykowy. 

Kolejne kierunki rozwoju systemu obejmują integrację z~innymi technologiami, takimi jak sensory LiDAR lub kamery RGB-D, jednakże ich użyteczność w~kontekście pracy wymaga dalszych badań. Należy również rozwinąć moduł aplikacji klienckiej dla systemu iOS, aby umożliwić korzystanie z~systemu jak największej liczbie użytkowników. Możliwe również, że w~przyszłości, inferencja będzie mogła odbywać się lokalnie na urządzeniu mobilnym, co pozwoliłoby na zwiększenie prywatności użytkowników i~odejście od konieczności stosowania jednostki serwerowej. 

Wymienione kierunki rozwoju systemu stanowią jedynie propozycje, które mogą przyczynić się do zwiększenia jego skuteczności i~użyteczności. Nie są to wszystkie możliwe kierunki, a~jedynie te które zostały zidentyfikowane w~toku realizacji pracy i~mogą przynieść wymierne korzyści w~przyszłości. Warto również podkreślić, że dalsze badania i~rozwój powinny odbywać się w~ścisłej współpracy z~ośrodkami zajmującymi się wsparciem osób niewidomych i~słabowidzących, aby zapewnić, że system będzie odpowiadał na rzeczywiste potrzeby użytkowników końcowych.